% Root supplied the cutoff outline; /root/p081_weighted_repair_fresh
% supplied this full source-support proof. Independent review is pending.
\begin{proof}
\emph{Actual domains and finite coefficients.}
In the coordinates \(t,q\), the change \(w=t+q\) identifies
\(E\) with the domain
\[
             0<v(t)<\delta,\qquad v(q)>v(t)
\]
in \((\mathbf G_{m,L}^2)^{\mathrm{an}}\).
Indeed \(|q|<|t|\) implies \(|w|=|t|\), so this domain
lies in \(D\); its inverse is \(q=w-t\).
These are analytic maps of the actual \(F\)-spaces, not
base changes of the current. Equivalently, the coordinates
\(t,q/t\) identify \(E\) with a standard annulus times a
punctured open unit disc over \(L\).

For either displayed moment, the inverse image of a compact
subset of its target is compact. In the coordinates \(t,w\),
respectively \(t,q\), it is a closed subset of a compact
closed polyannulus whose radii stay away from all excluded
boundaries. Arbitrary positive real radii are allowed; no
radius is required to be the absolute value of a field
element. Viewing the same polyannulus over \(F\) does not
change its topology or compactness.

We recall the short coefficient argument to avoid assuming
a source-calibration representation. For either target \(Q\),
pairing \(T\) with the pullback of \(H\Omega_{xy}\) gives a
positive distribution \(\nu\) on \(Q\). For each compact
\(K\subset Q\), choose a nonnegative compactly supported
smooth cutoff \(\theta\) equal to one on \(K\). Positivity
gives
\[
 |\nu(H)|\leq\|H\|_\infty\nu(\theta)
       \quad\hbox{if }\operatorname{supp}H\subset K.
\]
Thus \(\nu\) is a locally finite positive Radon measure.
Deleting \(\mathrm d'x\) or \(\mathrm d'y\) from the target
top form, with the appropriate sign, gives compactly
supported \((1,2)\)-form primitives. Their pullbacks have
compact analytic support by the preceding paragraph.
Since \(\mathrm d'T=0\), both distributional partial
derivatives of \(\nu\) vanish.

On a ball compactly contained in \(Q\), convolve with a
Euclidean mollifier of sufficiently small support.
The resulting smooth function has zero first derivatives
and hence is constant. Testing against a smooth function
of integral one shows that these constants have a finite
limit as the mollifier shrinks. Thus \(\nu\) is a finite
constant multiple of Lebesgue measure on each such ball.
Overlaps and connectedness of \(Q_D,Q_E\) give the unique
nonnegative constants in
\eqref{eq:p213-mixed-old} and \eqref{eq:p213-mixed-new}.
This is a distribution calculation on the targets, not
a slice of \(T\). The proper-support point agrees with
\cite[Section 7.1.11, p.~159]{CLD25}.

\emph{The three ordinary source regions.}
Let \(U=D\setminus\{q=0\}\). On this ordinary open set,
the functions
\[
 r=v(t),\qquad y=s-r=v(w/t),\qquad
 \zeta=v(q/t)
\]
are ordinary smooth logarithms of analytic units.
Their \(\mathrm d'\mathrm d''\) derivatives vanish;
see \cite[Section 18.1.2, p.~501]{CLD25}.
The strict ultrametric inequality gives the following
identities on whole ordinary open regions:
\begin{equation}
 \label{eq:p213-mixed-star}
 \begin{array}{c|c}
 y<0 & \zeta=y,\\
 y>0 & \zeta=0,\\
 \zeta>0 & y=0 .
 \end{array}
\end{equation}
For example \(y<0\) means \(|w/t|>1\), so
\(|w/t-1|=|w/t|\). If \(y>0\), then
\(|w/t|<1\) and \(|w/t-1|=1\).
If \(y=0\), the triangle inequality gives
\(|w/t-1|\leq1\), hence \(\zeta\geq0\).
Consequently these regions together with \(y=\zeta=0\)
cover \(U\). The last central set is not asserted to
be an ordinary-open constancy set for either function.

Choose \(0<a<b<\delta\) and
\(\varphi\in C_c^\infty((a,b))\) with
\(\int\varphi(r)\,\mathrm dr=1\). Choose numbers
\[
          0<\epsilon<\kappa<
                      \min(a,\delta-b).
\]
Choose \(\eta\in C_c^\infty((-\kappa,\kappa))\)
equal to one on a neighbourhood of
\([-\epsilon,\epsilon]\), and choose
\(\chi\in C_c^\infty((-\epsilon,\epsilon))\)
equal to one on a neighbourhood of zero.
On \(U\) define the ordinary smooth \((1,2)\)-form
\[
 \beta=\varphi(r)\eta(y)\chi(\zeta)\,
           \mathrm d'r\wedge\mathrm d''r\wedge\mathrm d''\zeta.
\]

\emph{Compact support on the entire analytic source.}
The possible support of \(\beta\) satisfies
\[
 a\leq r\leq b,\qquad |y|\leq\kappa,\qquad
 |\zeta|\leq\epsilon .
\]
Since \(s=r+y\), the first two inequalities put it in the
compact closed polyannulus
\[
 P=\{a\leq v(t)\leq b,\quad
                a-\kappa\leq v(w)\leq b+\kappa\}
       \Subset D.
\]
The third inequality also bounds
\[
 |q|=|t|\,|q/t|\ \geq\ \mathrm e^{-(b+\epsilon)}>0.
\]
The subset of \(P\) satisfying
\[
       \mathrm e^{-\epsilon}\leq |q/t|
                                     \leq\mathrm e^\epsilon
\]
is closed and compact and is contained in \(U\).
The support of \(\beta\) is a closed subset of this compact
set. Thus it is compactly contained in \(U\).
In particular \(\beta\) vanishes near the excluded divisor
\(q=0\) and extends by zero to a smooth compactly supported
form on \(D\). This checks support on the entire source,
not just on any chosen Gauss plane.

\emph{Cancellation before pairing with the current.}
The derivative of \(\varphi(r)\) in
\(\mathrm d'\beta\) vanishes by repetition of
\(\mathrm d'r\). The possible cutoff error from \(\eta\)
is
\[
 \varphi(r)\eta'(y)\chi(\zeta)\,
 \mathrm d'y\wedge\mathrm d'r\wedge
                         \mathrm d''r\wedge\mathrm d''\zeta .
\]
It vanishes as a source form. On \(y<0\) one has
\(\zeta=y\), while \(\eta'(y)\ne0\) forces
\(|y|>\epsilon\), where \(\chi(y)=0\).
On \(y>0\) one has \(\mathrm d''\zeta=0\) on
the entire ordinary open region. At every point with
\(y=0\), the factor \(\eta'(y)\) vanishes on an
ordinary neighbourhood, by continuity and the choice
of \(\eta\). Therefore
\[
 \mathrm d'\beta
   =\varphi(r)\eta(y)\chi'(\zeta)\,
       \mathrm d'r\wedge\mathrm d''r\wedge
                       \mathrm d'\zeta\wedge\mathrm d''\zeta.
\]
The sign here is positive: moving \(\mathrm d'\zeta\)
past \(\mathrm d'r,\mathrm d''r\) involves two
interchanges.

There is an ordinary neighbourhood of every point
with \(\zeta=0\) on which this derivative vanishes,
because \(\chi'\) is zero on an interval around zero.
Together with the cancellation just proved, this
excludes a central-face contribution without any
hypothesis on the current's lower-dimensional support.
The derivative is the sum of two compactly supported
forms on the disjoint regions \(\zeta<0\) and
\(\zeta>0\), each extended by zero. Their supports
stay a positive distance from \(\zeta=0\) in its
real coordinate.

\emph{The old and new integrals.}
On \(\zeta<0\), \eqref{eq:p213-mixed-star} gives
\(y=\zeta=s-r\), and \(\eta(y)=1\) on the support
of \(\chi'(y)\). The coordinate change
\((r,s)\mapsto(r,s-r)\) is an integral shear of
determinant one. In the top wedge, all terms involving
a repeated \(\mathrm d'r\) or \(\mathrm d''r\)
vanish. The negative piece is therefore the old
moment pullback with coefficient
\[
       \varphi(r)\chi'(s-r)\mathbf 1_{s<r}.
\]
This is a smooth compact target test: \(\chi'\)
vanishes near zero, and \(\epsilon<\min(a,\delta-b)\)
keeps its support inside \(Q_D\). For each relevant
\(r\), the interval \((-r,0)\) contains the negative
support of \(\chi'\). By
\eqref{eq:p213-mixed-old}, its pairing is
\[
 \begin{split}
 I_-&=B_T\int_a^b\varphi(r)
       \left(\int_0^r\chi'(s-r)\,\mathrm ds\right)\mathrm dr\\
    &=B_T\int_a^b\varphi(r)
               \bigl(\chi(0)-\chi(-r)\bigr)\,\mathrm dr
      =B_T .
 \end{split}
\]

On \(\zeta>0\), one has \(y=0\), \(\eta(y)=1\),
and \(h=r+\zeta>r\). Since \(\varphi\) is supported
in \((a,b)\), this piece lies in \(E\).
The change \((r,h)\mapsto(r,h-r)\) is again an
integral shear of determinant one. The positive
piece is the new moment pullback with coefficient
\(\varphi(r)\chi'(h-r)\), a compactly supported
smooth target test on \(Q_E\). Consequently
\[
 \begin{split}
 I_+&=C_T\int_a^b\varphi(r)
       \left(\int_r^\infty\chi'(h-r)\,\mathrm dh\right)
                        \mathrm dr\\
    &=C_T\int_a^b\varphi(r)
                 \bigl(0-\chi(0)\bigr)\,\mathrm dr
      =-C_T .
 \end{split}
\]
Closedness on \(D\), applied to the compactly supported
zero extension of \(\beta\), now yields
\[
        0=\langle T,\mathrm d'\beta\rangle
                       =I_-+I_+=B_T-C_T.
\]
This proves \eqref{eq:p213-mixed-equality}.

\emph{Canonical factors and the original field.}
The coordinates \(t,w\) on \(D\), and \(t,q\) on \(E\),
are primitive coordinates on their actual torus
domains over \(L\). Their canonical top-dimensional
coefficients over \(L\) are one: the coordinate moments
have local degree one in the defining calibration
formula of \cite[Section 14.3.2, p.~407]{CLD25}.
When each same analytic space is viewed over \(F\),
its canonical calibration is multiplied by the full
degree \([L:F]\), by
\cite[Remark 14.4.3, p.~410]{CLD25}.
This includes ramified finite extensions. No separate
sheet count or additional lattice factor is introduced;
both coordinate shears above have determinant one.
The entire calculation for \(T\) was performed on
ordinary \(F\)-analytic domains with \(T\) restricted
as an \(F\)-current. There was no pullback, descent, or
reinterpretation of \(T\) as an \(L\)-current.
\end{proof}
